\section{Estado del arte}

El desarrollo de juegos de rol multijugador en línea (MMO RPG) ha sido ampliamente explorado por títulos como World of Warcraft,
RuneScape y Final Fantasy XIV, los cuales sentaron las bases de los mundos persistentes con miles de jugadores conectados simultáneamente.
Sin embargo, estas propuestas se centran en infraestructuras cerradas que no permiten a los usuarios diseñar sus propios entornos,
limitando la personalización y la capacidad de experimentación dentro del ecosistema de juego. 
\vspace{0.25cm}

Por otro lado, motores como RPG Maker, Unity y Unreal Engine han democratizado la creación de videojuegos,
aunque su alcance técnico se orienta más hacia experiencias de un solo jugador o de multijugador reducido,
sin ofrecer de forma nativa soluciones que soportan un alto volumen de transacciones y conexiones en tiempo real.
\vspace{0.25cm}

En los últimos años, plataformas como Roblox y Minecraft han marcado un punto de inflexión al integrar la creación de mundos por parte de los jugadores
con experiencias multijugador de gran escala, convirtiéndose en referentes del contenido generado por usuarios (UGC). 
\vspace{0.25cm}

A nivel tecnológico, el avance de infraestructuras en la nube, como \textit{AWS GameLift}, \textit{Microsoft PlayFab} o \textit{Photon}, ha permitido que los desarrolladores 
gestionen de manera más eficiente la escalabilidad, el balanceo de carga y la persistencia de datos en tiempo real.
\vspace{0.25cm}

No obstante, aún existe un espacio de innovación en la intersección entre la accesibilidad de herramientas de creación como RPG Maker 
y la robustez de arquitecturas masivas propias de los MMO RPG, lo cual abre la oportunidad para proyectos que busquen combinar ambas dimensiones en una misma propuesta.
